%
% NOTEBOOKS
%

% #1 Notebook Name
% #2 Cell Name
% For use in sentence
\newcommand{\cellref}[2]{\hyperref[cell:#1-#2]{$#1_{\ref{cell:#1-#2}}$}}
\newcommand{\nbref}[2]{cell \hyperref[cell:#1-#2]{$#1_{\ref{cell:#1-#2}}$}}
% For use inside math
\WithSuffix\newcommand\nbref*[2]{\hyperref[cell:#1-#2]{#1_{\ref{cell:#1-#2}}}}
% For use at start of sentence
\newcommand{\Nbref}[2]{Cell \hyperref[cell:#1-#2]{$#1_{\ref{cell:#1-#2}}$}}

\newcommand{\rnbref}[1]{\hyperref[nb:#1]{#1}}


% #1 - Nodebook Name
% #2 - Cell Name
\newcommand{\nbCellLabel}[2]{cell:#1-#2}

% #1 - Cell Name
% #2 - Cell Code
\newcommand{\cell}[2]{
    \edef\lastcellid{\thecellid}
    \refstepcounter{cellid}
    \label{\nbCellLabel{\nbName}{#1}}

    \node[below=1mm of \lastcellid, inner sep=2mm, text width=\nbWidth, rounded corners, thick, draw=blue!20, fill=blue!10] (\thecellid) {\small\begin{obeylines}#2\end{obeylines}};
    \node[below left=0.5mm and 0mm of \thecellid.north east, @link] {\footnotesize $\nbName_\thecellid$};
}

\newcounter{notebookId}
\setcounter{notebookId}{0}
\newcounter{cellid}
\counterwithin*{cellid}{notebookId}

% #1 - Width
% #2- Notebook Name (used in cell label)
\newenvironment{notebook}[2][4cm]{
    \edef\nbName{#2}
    \edef\nbWidth{#1}
    \refstepcounter{notebookId}\label{nb:\nbName}

    \begin{center}
    \begin{tikzpicture}
    \node[color=blue!40] (0) {\small Notebook \nbName};
}{
    \begin{pgfonlayer}{background}
        \draw[draw=blue!30, fill=blue!1, rounded corners, thick]
        ($(current bounding box.south west) - (4pt, 4pt)$)
        rectangle
        ($(current bounding box.north east) + (4pt, 2pt)$);
    \end{pgfonlayer}
    \end{tikzpicture}
    \end{center}
}
